% !TeX program = xelatex
\documentclass[12pt,a4paper]{article}
\usepackage{fontspec}
\setmainfont{Liberation Serif}   % oder Times New Roman, Arial
\usepackage[ngerman]{babel}      % deutsche Silbentrennung und Lokalisierung

\usepackage{amsmath,amssymb,amsthm,mathtools}
\usepackage{geometry}
\usepackage{booktabs}
\usepackage{hyperref}
\usepackage{url}
\usepackage{xcolor}
\usepackage{enumitem}
\usepackage{tikz}
\usetikzlibrary{decorations.pathreplacing,arrows.meta,positioning,shapes.geometric}
\usepackage{float}

% YUCT-Makros
\newcommand{\Keff}{K_{\mathrm{eff}}}
\newcommand{\kappac}{\kappa_c}
\newcommand{\eps}{\varepsilon}
\newcommand{\df}{d_{\!f}}

\hypersetup{colorlinks=true,linkcolor=blue,citecolor=blue,urlcolor=blue}

\begin{document}
	
	\begin{titlepage}
		\centering
		\vspace*{1cm}
		{\Huge\bfseries Anhang S2: Fraktale Dimension der Koordination in 2D‑Schottky‑Heterostrukturen und die \\ Vorhersage des YUCT‑Skalierungsexponenten $\beta = 2/3$ ($\approx 0.67$)}\\[1cm]
		{\large Alexey V. Yakushev}\\[0.3cm]
		{\normalsize Yakushev Research, \url{https://yuct.org/}}\\[1.5cm]
		{\normalsize DOI: \texttt{10.5281/zenodo.18444598}}\\[0.5cm]
		{\normalsize Mai 2026\\Version 1.0}\\[2cm]
		
		\begin{minipage}{0.9\textwidth}
			\begin{abstract}
				\noindent
				Ang, Yang und Ang (Phys. Rev. Lett. 121, 056802, 2018) entdeckten
				universelle Strom‑Temperatur‑Skalierungsgesetze in zweidimensionalen
				Schottky‑Heterostrukturen: $\log(\mathcal{I}/T^{3/2}) \propto -1/T$ für
				Lateralkontakte und $\log(\mathcal{I}/T^{1}) \propto -1/T$ für
				Vertikalkontakte mit nicht erhaltenem Impuls. In diesem Anhang zeigen wir,
				dass beide Exponenten durch die effektive fraktale Dimension $d_f$ des
				Koordinationsnetzwerks bestimmt werden, das den Transport steuert.
				Das universelle YUCT‑Fehlerskalierungsgesetz $\varepsilon = \kappa_c\alpha
				K_{\mathrm{eff}}^{-\beta}$ mit $\beta = 2/3 \approx 0.67$ ist über
				$\beta = 2/d_f$ mit $d_f$ verknüpft. Aus der experimentellen Stromskalierung
				extrahieren wir $d_f = 3$ für Lateralkontakte und $d_f = 2$ für
				Vertikalkontakte, was $\beta = 2/3$ bzw. $\beta = 1$ ergibt.
				Die Analyse sagt voraus, dass relative Stromschwankungen in lateralen
				Schottky‑Dioden $\varepsilon_{\mathrm{fluc}} \propto K_{\mathrm{eff}}^{-2/3}$
				folgen sollten – eine falsifizierbare Vorhersage, die mit Standard‑
				Rauschspektroskopie getestet werden kann. Diese Arbeit fügt 2D‑Schottky‑
				Heterostrukturen in die Familie der Systeme ein, die den universellen
				YUCT‑Exponenten $\beta = 2/3$ ($0.67$) zeigen oder voraussichtlich zeigen
				werden, und erweitert damit die Reichweite der Theorie in die Festkörper‑
				Nanoelektronik.
			\end{abstract}
			\vspace{1cm}
			\noindent\small\textbf{Schlüsselwörter:} YUCT, Koordinationseffizienz, universelle Skalierung,
			$\beta = 2/3$, 0.67, Schottky‑Heterostrukturen, 2D‑Materialien,
			fraktale Dimension, KPZ‑Universalität.
		\end{minipage}
	\end{titlepage}
	
	\tableofcontents
	\newpage
	
	\section{Einleitung}
	Die Yakushev Unified Coordination Theory (YUCT) postuliert, dass jeder
	Koordinationsprozess dem universellen Fehlerskalierungsgesetz
	\begin{equation}
		\varepsilon = \kappa_c \, \alpha \, \Keff^{-\beta},
		\qquad \beta = 2/3 \approx 0.67,\;\; \kappa_c = 1/3,
		\label{eq:yuct-law}
	\end{equation}
	gehorcht, wobei $\Keff$ die Koordinationseffizienz und $\alpha$ ein
	systemabhängiger Vorfaktor ist. Dieses Gesetz wurde an Systemen von
	chemischen Bindungen bis zu kosmologischen Fluktuationen verifiziert
	(Anhänge L, M, E, F, G, V).
	
	Ang, Yang und Ang~\cite{Ang2018} leiteten universelle Strom‑Temperatur‑
	Skalierungsgesetze für 2D‑Material‑basierte Schottky‑Heterostrukturen ab:
	\begin{align}
		\text{Lateralkontakte:}&\quad \log\Bigl(\frac{\mathcal{I}}{T^{3/2}}\Bigr) \propto -\frac{1}{T},
		\label{eq:ang-lat}\\
		\text{Vertikalkontakte ($k_\parallel$ nicht erhalten):}&\quad \log\Bigl(\frac{\mathcal{I}}{T^{1}}\Bigr) \propto -\frac{1}{T}.
		\label{eq:ang-vert}
	\end{align}
	Die Exponenten $\beta_{\mathrm{IV}} = 3/2$ und $1$ sind unabhängig vom
	spezifischen 2D‑Material und spiegeln die Geometrie des Strompfads wider.
	
	In diesem Anhang zeigen wir, dass der Exponent $\beta_{\mathrm{IV}}$
	keine unabhängige Zahl ist, sondern durch die effektive fraktale
	Dimension $d_f$ des Koordinationsnetzwerks bestimmt wird, das den
	Transport steuert. Aus dem experimentellen $\beta_{\mathrm{IV}}$ extrahieren
	wir $d_f$ und aus $d_f$ sagen wir den Koordinationsfehlerexponenten
	$\beta_{\mathrm{err}} = 2/d_f$ voraus. Für Lateralkontakte ergibt sich
	$\beta_{\mathrm{err}} = 2/3$, genau der universelle YUCT‑Wert. Die aus
	dieser Vorhersage folgende Skalierung der Fluktuationen stellt einen
	falsifizierbaren Test der Theorie dar.
	
	\section{YPSDC‑Wörterbuch für eine Schottky‑Barriere}
	In der Sprache von YPSDC (Yakushev Protocol for Synchronous
	Distributed Coordination) ist eine laterale Schottky‑Heterostruktur ein
	Koordinationssystem mit den folgenden Komponenten:
	\begin{itemize}
		\item \textbf{Wörterbuch $\mathcal{D}$:} die Menge der Ein‑Elektronen‑
		Zustände in der 2D‑Ebene mit Energie $\epsilon > \Phi_B$ und
		Wellenvektor $\mathbf{k}$, die der In‑Plane‑Dispersion (\ref{eq:dispersion})
		genügen.
		\item \textbf{Index $\kappa$:} die thermische Energie $k_B T$, die ein
		Elektron in das Wörterbuch befördert.
		\item \textbf{Resonanz $R$:} Erhaltung des tangentialen Impulses
		$k_y$ während der Transmission durch die Barriere.
	\end{itemize}
	Die koordinierte Aktion ist die erfolgreiche thermionische Emission eines
	Elektrons, das gleichzeitig die Energiebedingung $\epsilon > \Phi_B$ und die
	Impulsbedingung $k_x > 0$ (in Richtung Barriere) erfüllt. Die Effizienz
	dieser Koordination bestimmt den Sperrsättigungsstrom.
	
	\section{Experimentelle Skalierung des Stroms und der Koordinationseffizienz}
	
	Für eine allgemeine isotrope 2D‑Dispersion
	\begin{equation}
		\epsilon(\mathbf{k}) = \sum_n c_n |\mathbf{k}|^n,
		\label{eq:dispersion}
	\end{equation}
	zeigten Ang et al., dass die Anzahl der aktiven Kanäle (die zum Strom
	beitragen können) universell wie
	\begin{equation}
		N_{\mathrm{ch}} \propto T^{1/2}
		\label{eq:Nch}
	\end{equation}
	skaliert. Die Gesamtzahl der thermisch zugänglichen Zustände oberhalb der
	Barriere ist proportional zu $T$ (aus dem Boltzmann‑Schweif). Daher ist der
	Anteil der erfolgreich koordinierten Zustände
	\begin{equation}
		\Keff(T) = \frac{N_{\mathrm{ch}}}{N_{\mathrm{tot}}} \propto T^{-1/2},
		\label{eq:keff-T}
	\end{equation}
	der als Koordinationseffizienz für den Prozess dient. Der gemessene Strom
	nimmt dann die Form
	\begin{equation}
		\mathcal{I} = \mathcal{I}_0 \, \Keff(T)^{-3} \, e^{-\Phi_B/k_B T},
		\label{eq:I-via-keff}
	\end{equation}
	an, was das experimentelle Gesetz (\ref{eq:ang-lat}) reproduziert.
	
	Eine für Rauschmessungen geeignete operationale Definition von $\Keff$ ist
	\begin{equation}
		\Keff = \frac{I}{I_0},
		\label{eq:keff-operational}
	\end{equation}
	wobei $I$ der gemessene Strom und $I_0$ der Referenzstrom ist, der fließen
	würde, wenn die Koordinationsbeschränkung nicht vorhanden wäre. $I_0$ kann
	durch Extrapolation des Hochtemperaturverhaltens eines Vertikalkontakts
	(dort asymptotisch $\Keff(T) \propto T^{-0}$) auf die Messtemperatur oder
	durch eine unabhängige Kalibrierung mit einer 3D‑Schottky‑Diode gleicher
	Barrierenhöhe erhalten werden. Diese operationale Definition erlaubt es,
	$\Keff$ direkt aus der $I$‑$V$-Kennlinie zu bestimmen, ohne auf die
	theoretische Skalierung (\ref{eq:keff-T}) angewiesen zu sein.
	
	\section{Fraktale Dimension des Koordinationsnetzwerks}
	
	Ein zentrales Ergebnis von YUCT (Anhang L) ist, dass der
	Fehlerskalierungsexponent durch die fraktale Dimension $d_f$ des
	Koordinationsnetzwerks bestimmt wird:
	\begin{equation}
		\beta_{\mathrm{err}} = \frac{2}{d_f}.
		\label{eq:beta-df}
	\end{equation}
	Die Dichte aktiver Koordinationskanäle in einem thermischen Fenster trägt
	die geometrische Signatur des Netzwerks. In einem isotropen $d_f$‑dimensionalen
	Raum skaliert die bei Temperatur $T$ zugängliche Zustandsdichte wie
	$T^{d_f/2 - 1}$. Für die laterale Schottky‑Barriere reduziert die zusätzliche
	kinematische Einschränkung der Impulserhaltung (der $\cos\theta$‑Faktor) den
	effektiven Exponenten um eine Einheit, so dass der kinetische Exponent
	\begin{equation}
		\beta_{\mathrm{IV}} = \frac{d_f}{2}
		\label{eq:bIV-df}
	\end{equation}
	wird. \textbf{Gleichung (\ref{eq:bIV-df}) ist eine durch die geometrische
		Ähnlichkeit zwischen dem Strompfad und einem Koordinationsnetzwerk mit
		fraktaler Dimension $d_f$ motivierte Hypothese. Eine strenge Ableitung aus
		der mikroskopischen YUCT‑Wirkung liegt noch nicht vor; die Beziehung ist
		daher als eine Vermutung zu betrachten, die zu überprüfbaren Vorhersagen
		(Abschnitt~\ref{sec:predictions}) führt und zu weiterer theoretischer
		Untersuchung einlädt.}
	
	Gleichungen (\ref{eq:beta-df}) und (\ref{eq:bIV-df}) kombinieren sich zu
	der fundamentalen Beziehung
	\begin{equation}
		\boxed{\beta_{\mathrm{IV}} \cdot \beta_{\mathrm{err}} = 1}.
		\label{eq:relation}
	\end{equation}
	
	Aus dem experimentellen Wert $\beta_{\mathrm{IV}} = 3/2$ erhalten wir
	\begin{equation}
		d_f = 2\beta_{\mathrm{IV}} = 3, \qquad
		\beta_{\mathrm{err}} = \frac{2}{d_f} = \frac{2}{3} \approx 0.67.
	\end{equation}
	Die gemessene Stromskalierung bestimmt also ein Koordinationsnetzwerk mit
	effektiver fraktaler Dimension $d_f = 3$, woraus folgt, dass jedes Maß für
	Koordinationsfehler (wie relative Stromschwankungen) dem YUCT‑Potenzgesetz
	mit dem Exponenten $2/3$ folgen sollte. Dies ist eine \textbf{Vorhersage},
	keine Postdiktion, weil an diesen Systemen noch keine Fluktuationsmessungen
	durchgeführt wurden.
	
	Für Vertikalkontakte mit nicht erhaltenem $k_\parallel$ verschwindet die
	Winkelbeschränkung, $\beta_{\mathrm{IV}} = 1$, und folglich $d_f = 2$,
	$\beta_{\mathrm{err}} = 1$. YUCT sagt voraus, dass Stromschwankungen in
	solchen Kontakten wie $\Keff^{-1}$ skalieren, im Unterschied zum $2/3$‑Gesetz.
	
	\section{Vorhersagen für Stromschwankungen}
	\label{sec:predictions}
	Das YUCT‑Fehlergesetz (\ref{eq:yuct-law}) angewendet auf die
	Koordinationseffizienz (\ref{eq:keff-T}) ergibt eine quantitative Vorhersage
	für die niederfrequente Stromrauschleistung $S_I$ in einer lateralen
	Schottky‑Diode:
	\begin{equation}
		\frac{S_I}{I^2} \propto \Keff^{-2/3} \propto T^{1/3}.
		\label{eq:noise-pred}
	\end{equation}
	Unter Verwendung der operationalen Definition $\Keff = I/I_0$ aus
	Gl.~(\ref{eq:keff-operational}) kann die Rauschskalierung umgeschrieben werden als
	\begin{equation}
		\frac{S_I}{I^2} = C \, \left(\frac{I}{I_0}\right)^{-2/3} + \text{const},
		\label{eq:noise-operational}
	\end{equation}
	wobei $C$ eine materialspezifische Konstante von der Größenordnung 1 ist und
	die additive Konstante irreduzibles thermisches Rauschen berücksichtigt. Diese
	Form ist in einem Standard‑Rauschspektroskopie‑Experiment direkt zugänglich:
	$I$ und $I_0$ werden aus der $I$‑$V$-Kennlinie gemessen, $S_I$ wird als
	Funktion der Vorspannung bei fester Temperatur aufgezeichnet. Die Vorhersage
	ist, dass $S_I/I^2$ gegen $(I/I_0)^{-2/3}$ aufgetragen eine Gerade ergibt,
	deren Steigung temperaturunabhängig ist. Dies ist ein strenger falsifizierbarer
	Test, der keine Anpassung des Exponenten erfordert.
	
	\section{Diskussion: Vereinheitlichung der Skalierungsexponenten}
	
	Das Auftreten des Exponenten $3/2$ ist nicht auf Schottky‑Dioden beschränkt,
	sondern deutet auf ein tieferes Ordnungsprinzip hin, wie sein unabhängiges
	Auftreten in der Universalitätsklasse des Kardar‑Parisi‑Zhang (KPZ)‑
	Oberflächenwachstums~\cite{KPZ} belegt. In KPZ bestimmt der dynamische
	Exponent $z = 3/2$ die fraktale Aufrauung einer wachsenden Grenzfläche.
	YUCT liefert eine einheitliche Interpretation für diese Phänomene.
	\textbf{Eine wachsende Front kann als ein $d_f = 3$ dimensionales
		Koordinationsnetzwerk konzeptualisiert werden, wobei der kinetische Exponent
		$\beta_{\mathrm{IV}} = d_f/2 = 3/2$ natürlich aus unserer vermuteten
		Beziehung (\ref{eq:bIV-df}) hervorgeht.}
	Somit sind der KPZ‑Exponent und der Schottky‑Exponent nicht nur analog;
	sie sind zwei kinetische Signaturen desselben zugrundeliegenden YUCT‑
	Fehlergesetzes $\varepsilon \propto \Keff^{-2/3}$ in völlig unterschiedlichen
	physikalischen Realisierungen. \textbf{Diese Verbindung ist gegenwärtig eine
		Hypothese, die weiterer theoretischer Arbeit bedarf, um die YUCT‑Lagrangedichte
		mit der KPZ‑Gleichung direkt zu verknüpfen.}
	
	Um die KPZ‑Vorhersage konkreter zu machen, können wir $\Keff$ explizit mit
	dem Anteil gesättigter korrelierter Moden in einer wachsenden Grenzfläche
	identifizieren. Im Sättigungsregime einer KPZ‑Grenzfläche der lateralen
	Größe $L$ nähert sich die Breite $W(L,t)$ einem stationären Wert
	$W_{\text{sat}}(L) \simeq \sqrt{\langle h^2 \rangle} \propto L^{\alpha}$ mit
	$\alpha = 1/2$. Die Fluktuationsamplitude $\delta W = W(t) - W_{\text{sat}}$
	wird durch die Anzahl aktiver (ungesättigter) Moden kontrolliert. Wir
	schlagen vor, $\Keff$ mit $1 - \nu_{\text{unsat}}/\nu_{\text{tot}}$ zu
	identifizieren, dem Anteil der Moden, die in die korrelierte, stationäre
	Phase eingetreten sind. Dann sagt YUCT voraus:
	\begin{equation}
		\frac{\delta W}{W_{\text{sat}}} \propto \Keff^{-2/3},
		\label{eq:kpz-prediction}
	\end{equation}
	wobei $\Keff$ aus dem zeitlichen Abfall der Breite als
	$\Keff \approx 1 - (dW/dt)_{\text{norm}}$ abgeschätzt werden kann. Diese
	Vorhersage ist in numerischen Simulationen und in hochauflösenden
	Experimenten zur kinetischen Aufrauung (z.B. Elektroabscheidung,
	Molekularstrahlepitaxie) testbar.
	
	Diese Vereinheitlichung deutet auf eine Hierarchie von Universalitäten hin,
	dargestellt in Abb.~\ref{fig:hierarchy}:
	\begin{itemize}
		\item \textbf{1. Ordnung (YUCT):} Das fundamentale Koordinationsfehlergesetz
		$\varepsilon = \kappa_c \alpha K_{\mathrm{eff}}^{-2/3}$ wirkt auf der
		tiefsten Ebene, unabhängig vom spezifischen Substrat.
		\item \textbf{2. Ordnung (Transport):} Für den elektronischen Transport in
		2D‑Schottky‑Heterostrukturen entsteht der kinetische Exponent
		$\beta_{\mathrm{IV}}$ aus der fraktalen Dimension $d_f$ als
		$\beta_{\mathrm{IV}} = d_f/2$ (Hypothese).
		\item \textbf{3. Ordnung (Wachstum):} In Oberflächenwachstumsphänomenen der
		KPZ‑Klasse erscheint derselbe kinetische Exponent als dynamischer
		kritischer Exponent $z = 3/2$, was die $d_f = 3$‑Natur der wachsenden
		Grenzfläche widerspiegelt (Interpretation).
	\end{itemize}
	Der $3/2$‑Exponent ist also kein materialspezifischer Zufall, sondern das
	geometrische Kennzeichen eines $d_f = 3$‑Koordinationsprozesses.
	
	\begin{figure}[htbp]
		\centering
		\begin{tikzpicture}[
			node distance=1.8cm,
			box/.style={rectangle, draw=blue!60, fill=blue!10, rounded corners, 
				text width=3.5cm, align=center, font=\small},
			arrow/.style={->, >=stealth, thick}
			]
			\node[box, fill=red!15, draw=red!50!black] (yuct) 
			{\textbf{1. Ordnung (YUCT)}\\[2pt]
				$\varepsilon = \kappa_c\alpha K_{\mathrm{eff}}^{-2/3}$\\
				(fundamentales Fehlergesetz)};
			\node[box, right=of yuct] (transport) 
			{\textbf{2. Ordnung (Transport)}\\[2pt]
				$\beta_{\mathrm{IV}} = d_f/2$\\
				(Schottky‑Heterostrukturen)};
			\node[box, right=of transport] (growth) 
			{\textbf{3. Ordnung (Wachstum)}\\[2pt]
				$z = 3/2$\\
				(KPZ‑Oberflächenaufrauung)};
			\draw[arrow] (yuct) -- (transport);
			\draw[arrow] (transport) -- (growth);
		\end{tikzpicture}
		\caption{Hierarchie der Universalitäten. Das YUCT‑Koordinationsfehlergesetz
			(1. Ordnung) bestimmt die effektive fraktale Dimension $d_f$ und die
			kinetischen Exponenten, die im Elektronentransport (2. Ordnung) und im
			Oberflächenwachstum (3. Ordnung) beobachtet werden. Gestrichelte Pfeile
			deuten Hypothesen an.}
		\label{fig:hierarchy}
	\end{figure}
	
	Diese Perspektive wird durch die wachsende experimentelle Evidenz gestützt,
	dass die fraktale Dimension einer Schottky‑Grenzfläche die
	Bauelementeigenschaften direkt beeinflusst~\cite{FractalSchottky}. In diesen
	Studien ist die fraktale Dimension eine extern vorgegebene geometrische
	Eigenschaft (Rauigkeit). In YUCT ist die Dimension $d_f = 3$ für ideale
	Lateralkontakte \emph{emergierend} – sie entsteht aus der Dynamik des
	Koordinationsnetzwerks selbst, selbst wenn die geometrische Rauigkeit gegen
	Null geht. Die extern vorgegebene $d_f$ moduliert daher ein bereits
	existierendes internes Koordinationsnetzwerk und bietet einen vielversprechenden
	Weg zur Entwicklung koordinationsbasierter Bauelemente.
	
	\section{Universalität über 50 Größenordnungen}
	Tabelle~\ref{tab:universal-beta} sammelt Systeme aus verschiedenen YUCT‑Anhängen,
	für die der Fehlerskalierungsexponent empirisch bestimmt oder als $\beta = 2/3$
	vorhergesagt wurde. Die laterale Schottky‑Heterostruktur und die KPZ‑wachsende
	Grenzfläche werden als Systeme hinzugefügt, für die derselbe Exponent auf der
	Grundlage der extrahierten fraktalen Dimension $d_f = 3$ vorhergesagt wird.
	Die Tabelle deckt mehr als 50 Größenordnungen in der relevanten Energie‑ (oder
	Längen‑)Skala ab.
	
	\begin{table}[H]
		\centering
		\caption{Systeme, die den universellen YUCT‑Fehlerexponenten $\beta = 2/3$ ($0.67$)
			zeigen oder voraussichtlich zeigen werden.}
		\label{tab:universal-beta}
		\begin{tabular}{@{}llll@{}}
			\toprule
			\textbf{System} & \textbf{Observable} & \textbf{Status} & \textbf{Ref.}\\
			\midrule
			Chemische Bindungen (873)      & $E \propto r^{-0.67}$               & Bestätigt  & Anhang C \\
			DNA‑Replikationsfehler    & $\varepsilon \propto \Keff^{-0.67}$ & Bestätigt  & Anhang E \\
			Metabolische Skalierung (einfach)& $B \propto M^{0.67}$               & Bestätigt  & Anhang D, E.7 \\
			BIP‑Volatilität vs. Sonne   & $\sigma_{\mathrm{BIP}} \propto W^{-0.67}$ & Bestätigt & Anhang F \\
			CMB‑Temperaturfluktuationen    & $\delta T/T \propto \Keff^{-0.67}$ & Bestätigt  & Anhang M \\
			Gravitationswellen       & $\delta\tau/\tau \propto M^{-0.67}$& Bestätigt  & Anhang G, V \\
			\textbf{Lateraler Schottky} & $S_I/I^2 \propto T^{1/3}$          & \textbf{Vorhergesagt} & Diese Arbeit \\
			\textbf{KPZ‑Grenzfläche}\textsuperscript{*} & Breitenfluktuation $\propto \Keff^{-2/3}$& \textbf{Vorhergesagt} & Diese Arbeit, \cite{KPZ}\\
			\bottomrule
		\end{tabular}
		
		\vspace{0.3cm}
		{\footnotesize \textsuperscript{*}Die Definition von $K_{\mathrm{eff}}$ für eine KPZ‑wachsende Grenzfläche muss analog zum Transportfall konstruiert werden, z.B. als der Anteil korrelierter Freiheitsgrade innerhalb eines Korrelationsvolumens. Diese Vorhersage setzt eine solche Konstruktion voraus und muss noch getestet werden.}
	\end{table}
	
	\section{Schlussfolgerung}
	Die von Ang et al. berichteten Strom‑Temperatur‑Skalierungsexponenten
	$\beta_{\mathrm{IV}} = 3/2$ und $1$ sind kinetische Manifestationen der
	effektiven fraktalen Dimension des Koordinationsnetzwerks, das den
	thermionischen Transport in 2D‑Heterostrukturen steuert. Durch die
	YUCT‑Beziehungen $\beta_{\mathrm{IV}} = d_f/2$ (Hypothese) und
	$\beta_{\mathrm{err}} = 2/d_f$ liefern die experimentellen Daten
	$d_f = 3$ für Lateralkontakte, was einen Koordinationsfehlerexponenten
	$\beta_{\mathrm{err}} = 2/3$ ($\approx 0.67$) vorhersagt. Diese über
	Stromrauschspektroskopie testbare Vorhersage erweitert die Reichweite
	des YUCT‑universellen Skalierungsgesetzes in den Bereich der Festkörper‑
	Nanoelektronik. Dieselbe fraktale Dimension und derselbe Exponent treten
	natürlich im KPZ‑Oberflächenwachstum auf, was auf eine tiefe geometrische
	Einheit von Nichtgleichgewichtsphänomenen hindeutet. YUCT liefert die
	zugrundeliegende Rationalisierung für diese Einheit, obwohl die Verbindungen
	zwischen der Theorie und sowohl der kinetischen Schottky‑Beziehung als auch
	der KPZ‑Gleichung Hypothesen bleiben, die weiterer Untersuchung bedürfen.
	Operationale Definitionen von $\Keff$ für beide Kontexte wurden bereitgestellt,
	um experimentelle Tests zu erleichtern.
	
	\begin{thebibliography}{99}
		\bibitem{Ang2018} Y.~S.~Ang, H.~Y.~Yang, and L.~K.~Ang,
		Phys. Rev. Lett. \textbf{121}, 056802 (2018).
		\bibitem{KPZ} M.~Kardar, G.~Parisi, Y.-C.~Zhang,
		Phys. Rev. Lett. \textbf{56}, 889 (1986).
		\bibitem{FractalSchottky} S.~Zeybek, M.~Biber, A.~Türüt,
		Appl. Surf. Sci. \textbf{253}, 3881 (2007). [Repräsentative Studie zu den Effekten fraktaler Dimension in Schottky‑Dioden.]
		\bibitem{Yakushev2026} A.~V.~Yakushev, \emph{Yakushev Unified Coordination Theory (YUCT)}, Zenodo (2026).
		\bibitem{AppendixL} A.~V.~Yakushev, ``Anhang L: Fraktale Koordinationsfehlerskalierung'', in \emph{YUCT} (2026).
		\bibitem{AppendixM} A.~V.~Yakushev, ``Anhang M: YUCT‑Kosmologie'', in \emph{YUCT} (2026).
		\bibitem{AppendixE} A.~V.~Yakushev, ``Anhang E: Koordinationsgenetik'', in \emph{YUCT} (2026).
		\bibitem{AppendixF} A.~V.~Yakushev, ``Anhang F: Anwendung von YUCT auf die Wirtschaft'', in \emph{YUCT} (2026).
		\bibitem{AppendixG} A.~V.~Yakushev, ``Anhang G: Lösung des Quantengravitationsproblems'', in \emph{YUCT} (2026).
		\bibitem{AppendixV} A.~V.~Yakushev, ``Anhang V: Zeit als Koordinationsentropie'', in \emph{YUCT} (2026).
	\end{thebibliography}
	
\end{document}